% Estilo comum das figuras. A mesma cor e o mesmo marcador significam
% sempre a mesma série. Carregue este arquivo depois de pgfplots.
%
% Mapa fixo:
%   0% de maliciosos     azul, círculo, linha contínua
%   25% de maliciosos    verde, triângulo, linha contínua
%   50% de maliciosos    laranja, quadrado, linha tracejada
%   65% de maliciosos    vermelho, estrela, linha pontilhada
%   reputação maliciosa  o mesmo par da série de 65%
%   maioria simples      cinza-escuro, losango, linha contínua
%   reputação honesta    roxo, pentágono, linha contínua
%   EMA assimétrica      marrom, otimes, linha tracejada
%   Beta                 verde-azulado, oplus, linha traço-ponto
%
% A curva EMA do cenário de 65% reutiliza a série de 65%, porque é a mesma
% curva da figura que compara os percentuais.
%
% Largura e altura: redefina \figw e \figh antes de cada eixo.

\definecolor{serieZero}{RGB}{31,119,180}
\definecolor{serieVinte}{RGB}{44,160,44}
\definecolor{serieCinquenta}{RGB}{230,120,20}
\definecolor{serieSessenta}{RGB}{200,30,30}
\definecolor{serieMaioria}{RGB}{40,40,40}
\definecolor{serieHonesto}{RGB}{106,90,205}
\definecolor{serieAssimetrico}{RGB}{140,86,75}
\definecolor{serieBeta}{RGB}{0,128,128}
\definecolor{corTeste}{RGB}{130,130,130}

\newcommand{\figw}{0.92\textwidth}
\newcommand{\figh}{7.4cm}

\pgfplotsset{
  compat=1.18,
  tcc linha/.style={
    width=\figw,
    height=\figh,
    scale only axis,
    font=\normalsize,
    tick label style={font=\normalsize},
    label style={font=\normalsize},
    legend style={font=\normalsize},
    xlabel={Rodada},
    xmin=-1, xmax=53,
    ymin=-0.08, ymax=1.10,
    xtick={0,10,20,30,40,50},
    ytick={0,0.2,0.4,0.6,0.8,1},
    grid=major,
    grid style={black!12},
    clip=false,
    axis line style={line width=0.6pt},
    tick style={black, line width=0.6pt},
    legend cell align=left,
    legend columns=2,
    legend style={
      at={(0.5,-0.46)},
      anchor=north,
      draw=none,
      fill=none,
      font=\normalsize,
      /tikz/every even column/.append style={column sep=1.1em},
    },
    every axis plot/.append style={
      line width=1.15pt,
      opacity=0.6,
      mark size=2.6pt,
      mark repeat=10,
      % Índice 0-based: o ponto 9 é a rodada 10; depois, a cada 10 rodadas.
      mark phase=9,
      mark options={solid, opacity=1},
    },
  },
  serie zero/.style={
    color=serieZero, solid, mark=*,
    mark options={solid, fill=serieZero, opacity=1},
  },
  serie vinte/.style={
    color=serieVinte, solid, mark=triangle*,
    mark options={solid, fill=serieVinte, opacity=1},
  },
  serie cinquenta/.style={
    color=serieCinquenta, dashed, mark=square*,
    mark options={solid, fill=serieCinquenta, opacity=1},
  },
  serie sessenta/.style={
    color=serieSessenta, dotted, mark=star, mark size=3.5pt,
    mark options={solid, fill=serieSessenta, opacity=1},
  },
  serie malicioso/.style={serie sessenta},
  serie maioria/.style={
    color=serieMaioria, solid, mark=diamond*,
    mark options={solid, fill=serieMaioria, opacity=1},
  },
  serie honesto/.style={
    color=serieHonesto, solid, mark=pentagon*,
    mark options={solid, fill=serieHonesto, opacity=1},
  },
  serie ema assimetrica/.style={
    color=serieAssimetrico, dashed, mark=otimes, mark size=3pt,
    mark options={solid, opacity=1},
  },
  serie beta/.style={
    color=serieBeta, dashdotted, mark=oplus, mark size=3pt,
    mark options={solid, opacity=1},
  },
}

% Linha vertical em cada rodada de teste. Usar dentro do eixo.
\newcommand{\linhasDeTeste}{%
  \pgfplotsinvokeforeach{10,20,30,40,50}{%
    \draw[corTeste, dashed, line width=0.8pt]
      ({rel axis cs:0,0} -| {axis cs:##1,0}) --
      ({rel axis cs:0,1} -| {axis cs:##1,0});
  }%
}

\newcommand{\legendaRodadaDeTeste}{%
  \addlegendimage{corTeste, dashed, line width=0.8pt}%
  \addlegendentry{Rodada de validação}%
}
